\documentclass{beamer}

\input{shared_header}

\begin{document}

%==================== SLIDE 1: TITLE ====================
{
\setbeamertemplate{footline}{}
\begin{frame}[noframenumbering]

\vspace{1.9cm}

\begin{center}
    {\fontsize{25}{29}\selectfont \color{tugreen}
    Case Study II:\\[0.25cm]
  Next-Lactation Milk Yield Forecasting}
\end{center}

\vspace{0.95cm}

\begin{flushleft}
{\large
Nixon M Paulson\\
Rohith B Chowdappa\\
Maria Cherian}
\end{flushleft}

\end{frame}
}

%==================== SLIDE 2: PROBLEM ====================
\begin{frame}{Problem Statement}

\setbeamercolor{item}{fg=tugreen}
\begin{itemize}
    \item Predict \textbf{total milk yield in the next lactation}
    \item Setup: features from lactation $n$ $\rightarrow$ target $= \texttt{mkg\_sum}(n{+}1)$
    \item Data: OptiKuh daily dairy records (1,714 cows, $\sim$3,400 lactations)
    \item Aggregation level: one row per cow (\texttt{lom}) + lactation (\texttt{lnr})
\end{itemize}

\end{frame}

%==================== SLIDE 3: INITIAL CHALLENGES ====================
\begin{frame}{Initial Challenges}

\setbeamercolor{item}{fg=tugreen}
\begin{itemize}
    \item Daily records must be aggregated without mixing lactations
    \item Disease fields are categorical text, not numeric flags
    \item Most cows appear in only one or two lactation episodes
\end{itemize}

\end{frame}

%==================== SLIDE 4: FEATURE ENGINEERING ====================
\begin{frame}{Feature Engineering Approach}

\setbeamercolor{item}{fg=tugreen}
\begin{itemize}
    \item Sort records by cow, lactation, and observation date
    \item Aggregate to lactation level with mean / min / max:
    \begin{itemize}
        \item Nutrition: \texttt{fts}, \texttt{ftsnel}, \texttt{ftsme}, \texttt{eb\_nel}, \texttt{eb\_me}
        \item Physiology: \texttt{bcs}, \texttt{gew}, blood markers (\texttt{bhb}, \texttt{nefa}, \texttt{glukose})
    \end{itemize}
    \item Milk signals: \texttt{mkg\_sum}, \texttt{mkg\_mean}, \texttt{mkg\_max}
    \item Target: \texttt{groupby(lom)["mkg\_sum"].shift(-1)}; drop last lactation per cow
    \item All model inputs forced to numeric (\texttt{float64}) before training
\end{itemize}

\end{frame}

%==================== SLIDE 5: DISEASE HANDLING ====================
\begin{frame}{Disease Handling}

\setbeamercolor{item}{fg=tugreen}
\textbf{Initial approach}
\begin{itemize}
    \item One-hot encoding of diagnosis text (26 sparse boolean columns)
    \item High dimensionality, low signal
\end{itemize}

\vspace{0.3cm}

\textbf{Improved approach}
\begin{itemize}
    \item \texttt{disease\_count}: total disease events per lactation
    \item \texttt{had\_disease}: binary flag (any event)
    \item \texttt{disease\_category\_count}: number of disease groups affected
\end{itemize}

\end{frame}

%==================== SLIDE 6: MODELS ====================
\begin{frame}{Models Used}

\setbeamercolor{item}{fg=tugreen}
\begin{itemize}
    \item \textbf{HistGradientBoostingRegressor}
    \item Handles missing values natively
    \item Hyperparameters: \texttt{max\_depth=8}, \texttt{learning\_rate=0.05}, \texttt{max\_iter=500}
    \item Baseline and improved pipelines use the same model for fair comparison
\end{itemize}

\end{frame}

%==================== SLIDE 7: KEY ISSUE ====================
\begin{frame}{Key Issue Discovered}

\setbeamercolor{item}{fg=tugreen}
\begin{itemize}
    \item \textbf{Missing temporal structure} was the main failure mode
    \item 74\% of cows contributed only one supervised row $\rightarrow$ lag features mostly empty
    \item 37\% of lactations had \texttt{mkg\_sum = 0} (incomplete milk recording window)
    \item Time-based split left early train folds without usable history
    \item Baseline excluded current-lactation \texttt{mkg\_sum} from features
\end{itemize}

\vspace{0.2cm}

\textbf{To note:} next-lactation yield needs cow-level milk history, not only within-lactation averages.

\end{frame}

%==================== SLIDE 8: IMPROVEMENTS ====================
\begin{frame}{Improvements Added}

\setbeamercolor{item}{fg=tugreen}
\begin{itemize}
    \item Lag features per cow (ordered by calving date):
    \begin{itemize}
        \item \texttt{prev\_lactation\_milk}, \texttt{prev\_prev\_lactation\_milk}
        \item \texttt{milk\_trend}, \texttt{cow\_avg\_milk} (prior lactations only)
    \end{itemize}
    \item Drop lactations with zero recorded milk (current or next)
    \item Simplified disease counts instead of one-hot explosion
    \item Remove all-NaN / constant columns before fitting
\end{itemize}

\end{frame}

%==================== SLIDE 9: RESULTS ====================
\begin{frame}{Results}

\small
\begin{center}
\begin{tabular}{l r r r}
\hline
 & \textbf{RMSE (kg)} & \textbf{MAE (kg)} & \textbf{R\textsuperscript{2}} \\
\hline
Baseline (before improvements) & 3,291 & 2,494 & 0.066 \\
Improved model                 & 2,556 & 1,941 & 0.453 \\
\hline
\end{tabular}
\end{center}

\vspace{0.4cm}

\setbeamercolor{item}{fg=tugreen}
\begin{itemize}
    \item 44 numeric features in final model
    \item Largest gain from lag features + data quality, not model complexity
\end{itemize}

\end{frame}

%==================== SLIDE 10: CONCLUSIONS ====================
\begin{frame}{Conclusions \& Next Steps}

\setbeamercolor{item}{fg=tugreen}
\textbf{Conclusions}
\begin{itemize}
    \item Correct temporal structure matters more than adding many raw features
    \item Simple sklearn baseline can reach useful performance with clean aggregation
\end{itemize}

\vspace{0.3cm}

\textbf{Next steps}
\begin{itemize}
    \item Tune features for cows with longer lactation histories
    \item Use designated time series models for better temporal modeling?
\end{itemize}

\end{frame}

\end{document}
